% Informe Project 01 - Computació Distribuïda
% Compila amb pdflatex (per exemple a Overleaf).
% ESTAT: esborrany parcial. Cobreix la part de la casa (sensors, gateway) i
% l'ingestor. Les seccions marcades com a PENDENT s'han d'omplir.

\documentclass[11pt,a4paper]{article}

\usepackage[utf8]{inputenc}
\usepackage[T1]{fontenc}
\usepackage[catalan]{babel}
\usepackage[margin=2.5cm]{geometry}
\usepackage{listings}
\usepackage{xcolor}
\usepackage{booktabs}
\usepackage{hyperref}

\lstset{
  basicstyle=\ttfamily\small,
  breaklines=true,
  frame=single,
  columns=fullflexible,
  keywordstyle=\color{blue!70!black},
  commentstyle=\color{green!40!black},
  showstringspaces=false
}

\newcommand{\pendent}[1]{\textcolor{red}{\textbf{[PENDENT: #1]}}}

\begin{document}

% ---------------------------------------------------------------- Portada
\begin{titlepage}
  \centering
  \vspace*{3cm}
  {\Large Universitat de Lleida\\[0.3cm] Grau en Enginyeria Informàtica\par}
  \vspace{1.5cm}
  {\large Computació Distribuïda\par}
  \vspace{1cm}
  {\Huge\bfseries Project 01\\[0.4cm] Plataforma IoT de confort per a llars\par}
  \vspace{2cm}
  {\large Agustí \pendent{cognoms}\\ Clara \pendent{cognoms}\par}
  \vfill
  {\large \pendent{data d'entrega}\par}
\end{titlepage}

\tableofcontents
\newpage

% ---------------------------------------------------------------- Objectiu
\section{Objectiu}

L'objectiu del projecte és simular una plataforma IoT que recull la
temperatura de les habitacions de diverses llars, la processa al núvol de
forma distribuïda i envia ordres a la bomba de calor de cada habitació per
mantenir el confort que vol cada usuari.

La simulació funciona a escala 1/60: un segon real equival a un minut simulat.
Cada lectura d'un sensor representa, per tant, un minut de la llar.

% ---------------------------------------------------------------- Arquitectura
\section{Arquitectura general}

El sistema té dues zones separades per xarxa:

\begin{itemize}
  \item \textbf{La llar} (una per usuari): els sensors de temperatura, un
        broker MQTT local (Mosquitto) i un \emph{gateway}.
  \item \textbf{El núvol}: un broker MQTT públic, l'\emph{ingestor}, Kafka,
        l'aplicació Kafka Streams (pipelines P2--P8), el servei
        d'emmagatzematge cap a InfluxDB i el servei \emph{commander}.
\end{itemize}

El flux d'una lectura és:
\begin{center}
\small
sensor $\rightarrow$ Mosquitto llar $\rightarrow$ gateway $\rightarrow$
Mosquitto núvol $\rightarrow$ ingestor $\rightarrow$ Kafka
(\texttt{raw-readings}) $\rightarrow$ emmagatzematge $\rightarrow$ InfluxDB
\end{center}
I el d'una ordre, el camí invers:
\begin{center}
\small
Kafka Streams (P4) $\rightarrow$ \texttt{heatpump-commands} $\rightarrow$
commander $\rightarrow$ Mosquitto núvol $\rightarrow$ gateway $\rightarrow$
Mosquitto llar $\rightarrow$ sensor
\end{center}

Cada llar es desplega amb el seu propi fitxer \emph{docker-compose} i el seu
propi \texttt{--project-name}. Totes les llars i el núvol comparteixen una
xarxa Docker externa, \texttt{global\_net}, que fa el paper d'Internet. Cada
llar té, a més, una xarxa privada (\texttt{internal\_net}) on només hi ha els
seus sensors, el seu broker i el seu gateway.

% ---------------------------------------------------------------- Components
\section{Descripció dels components}

\subsection{Sensor de temperatura}

El sensor ens ve donat i no es pot modificar. Cada segon publica una lectura
al topic configurat a \texttt{SENSOR\_TOPIC\_DATA} amb el format:
\begin{lstlisting}
{"timestamp": "2026-11-25T10:00:00.000000", "value": 21.4}
\end{lstlisting}
El \texttt{timestamp} és el temps \emph{simulat}, en format ISO~8601 i sense
zona horària. Quan els microsegons són zero, \texttt{isoformat()} no escriu la
part decimal, i el parseig de la P2 ho ha de tenir en compte. Cada lectura té
una probabilitat del 20\,\% de valer exactament 100 (aproximadament un 20\,\%
de les lectures): són errors del sensor, que la pipeline P2 ha de filtrar i
comptar.

El sensor escolta el topic \texttt{SENSOR\_TOPIC\_STATUS} i espera missatges
\texttt{\{"status": 1\}} (encendre) o \texttt{\{"status": 0\}} (apagar). El
seu codi no captura l'error si rep un missatge que no és JSON: l'excepció
atura el fil de xarxa del client MQTT, i el contenidor continua viu però deixa
de rebre ordres i d'enviar lectures. Per això el gateway només li reenvia
ordres que ja ha validat (secció~\ref{sec:gateway}).

\subsection{Llar: fitxers docker-compose}

Hem creat un fitxer per llar a partir del \emph{docker-compose} d'exemple:
\begin{itemize}
  \item \texttt{home1.yml}: usuari \texttt{user1}, habitacions
        \texttt{room1}, \texttt{room2} i \texttt{living\_room}.
  \item \texttt{home2.yml}: usuari \texttt{user2}, habitacions
        \texttt{room1} i \texttt{office}.
\end{itemize}
Les dues llars tenen habitacions diferents a propòsit, per demostrar que el
gateway no necessita cap configuració per habitació.

Els canvis respecte al fitxer d'exemple són:
\begin{itemize}
  \item El broker de la llar \textbf{no publica cap port} a la màquina
        amfitriona. Només hi parlen els contenidors de la mateixa llar, i així
        poden funcionar diverses llars a la mateixa màquina sense conflictes
        de port.
  \item Els sensors tenen \texttt{restart: on-failure}. El sensor es tanca si
        el broker encara no està llest quan arrenca, i \texttt{depends\_on}
        només espera que el contenidor del broker s'hagi iniciat, no que ja
        accepti connexions.
  \item Mantenim la imatge \texttt{eclipse-mosquitto:1.6.12} del fitxer
        d'exemple. Aquesta versió accepta connexions anònimes sense cap
        configuració; la 2.x necessitaria un fitxer de configuració.
  \item S'hi afegeix el servei \texttt{gateway}, connectat a les dues xarxes.
\end{itemize}

\subsection{Gateway}\label{sec:gateway}

El gateway (\texttt{gateway/gateway.py}, Python + \texttt{paho-mqtt}) fa de
pont entre el broker de la llar i el del núvol. Té dos clients MQTT, un per a
cada broker, i treballa en dues direccions:

\paragraph{Llar $\rightarrow$ núvol.}
Se subscriu a \texttt{+/temperature}. El comodí \texttt{+} captura totes les
habitacions, de manera que un sensor nou no requereix tocar el gateway. Per a
cada missatge:
\begin{enumerate}
  \item Extreu l'habitació del topic (\texttt{room1/temperature}
        $\rightarrow$ \texttt{room1}).
  \item Si el contingut és un objecte JSON vàlid, hi afegeix els camps
        \texttt{user} i \texttt{room}. Els camps originals del sensor
        (\texttt{timestamp} i \texttt{value}) no es modifiquen.
  \item Publica al núvol al topic \texttt{<user>/<room>/temperature}.
\end{enumerate}

\emph{Decisió: els missatges mal formats es reenvien, no es llencen.} La
pipeline P1 ha de guardar les dades ``exactament tal com les ha llegit el
sensor'', i P2 ha de comptar els errors de cada sensor. Si el gateway llencés
els missatges incorrectes, al núvol no hi arribarien mai i no es podrien
comptar. Com que un missatge mal format no es pot enriquir, es reenvia
intacte. La identitat (usuari i habitació) no es perd perquè també va dins
del topic.

\paragraph{Preferències de confort (llar $\rightarrow$ núvol).}
L'enunciat demana que el gateway també reculli les preferències de confort de
l'usuari. L'usuari les publica al broker de casa, al topic \texttt{settings}
(per exemple amb \texttt{mosquitto\_pub}), i el gateway les reenvia al núvol a
\texttt{<user>/settings}. El gateway no les valida: ho fa l'ingestor, que és
qui les escriu al topic compactat. Així la regla de validació és en un sol
lloc.

\paragraph{Núvol $\rightarrow$ llar.}
Se subscriu a \texttt{<user>/+/heatpump}, només les ordres del seu usuari, i
les publica al broker local a \texttt{<room>/heatpump}, que és el topic que
escolta el sensor. Abans de reenviar una ordre, comprova que tingui
exactament el format \texttt{\{"status": 0|1\}}. Així es protegeix el sensor,
que no tolera missatges incorrectes.

\paragraph{Robustesa.}
\begin{itemize}
  \item Si el broker encara no està disponible quan el gateway arrenca, ho
        torna a provar cada 5 segons.
  \item Les subscripcions es fan dins de \texttt{on\_connect}, de manera que
        es restauren automàticament després d'una reconnexió.
  \item Cada client té un \texttt{client\_id} únic
        (\texttt{gateway-<user>-cloud}). Dos clients amb el mateix
        identificador en un broker es desconnectarien l'un a l'altre.
  \item Cap al núvol es publica amb QoS~1 (entrega com a mínim una vegada).
        Això només cobreix el tram gateway $\rightarrow$ núvol: el sensor
        publica amb QoS~0, i al tram sensor $\rightarrow$ gateway no hi ha
        cap garantia d'entrega.
\end{itemize}

\subsection{Ingestor}

L'ingestor (\texttt{cloud/ingestor/ingestor.py}, Python + \texttt{paho-mqtt}
+ \texttt{confluent-kafka}) se subscriu al broker del núvol i publica cada
missatge a Kafka:

\begin{center}
\begin{tabular}{@{}lll@{}}
\toprule
Topic MQTT & Topic Kafka & Clau \\
\midrule
\texttt{<user>/<room>/temperature} & \texttt{raw-readings} & \texttt{<user>/<room>} \\
\texttt{<user>/settings} & \texttt{user-settings} & \texttt{<user>} \\
\bottomrule
\end{tabular}
\end{center}

\begin{itemize}
  \item Les lectures s'envien a \texttt{raw-readings} \textbf{sense
        interpretar-les}: el valor és el contingut exacte que arriba. Per
        això l'ingestor no pot fallar per culpa d'un missatge mal format,
        perquè no el descodifica.
  \item Els \emph{settings} sí que es validen (mode \texttt{comfort},
        \texttt{eco} o \texttt{away}; \texttt{min} i \texttt{max} opcionals,
        numèrics i amb $\mathit{min} < \mathit{max}$). El motiu és que
        \texttt{user-settings} és un topic compactat que P3 llegeix com a
        \texttt{KTable}: un valor incorrecte passaria a ser la configuració
        ``actual'' de l'usuari.
  \item El productor Kafka fa servir \texttt{acks=all} (la dada no es dona
        per enviada fins que el broker l'ha guardada) i
        \texttt{enable.idempotence=true} (els reintents no generen
        duplicats).
  \item \emph{Ampliacions respecte al temari.} Els dos punts següents no es
        van veure a classe. Els hem afegit per millorar la tolerància a
        fallades i l'escalabilitat de l'ingestor, que són requisits de
        l'enunciat (``consider efficiency, fault tolerance and performance'').
  \item El client MQTT fa servir una \textbf{sessió persistent}
        (\texttt{clean\_session=False}) amb un identificador estable. Mentre
        l'ingestor està aturat, el broker li guarda en cua els missatges
        QoS~1 i els hi entrega quan torna.
  \item Les subscripcions són \textbf{compartides}
        (\texttt{\$share/ingestors/...}, suportades per Mosquitto 1.6). Si
        s'arrenquen diversos ingestors, cada missatge el rep només un d'ells.
        Així l'ingestor es pot escalar horitzontalment sense duplicar dades a
        Kafka. Cada instància necessita un \texttt{client\_id} únic: per
        defecte es fa servir el nom de la màquina.
  \item En aturar el contenidor (\texttt{SIGTERM}), l'ingestor està pensat per
        tancar el client MQTT i fer \texttt{flush()} del productor, per no
        perdre missatges pendents. \pendent{verificar-ho a la prova de punta a
        punta.}
\end{itemize}

\subsection{Elecció de la clau de Kafka}

Kafka només garanteix l'ordre dels missatges \emph{dins d'una partició}, i la
clau decideix a quina partició va cada missatge. Fem servir
\texttt{<user>/<room>} com a clau de \texttt{raw-readings} perquè:
\begin{itemize}
  \item \textbf{Ordre per habitació:} totes les lectures d'una habitació van
        a la mateixa partició i en ordre. És el que necessiten les
        operacions amb estat per habitació: l'estat de la bomba de calor
        (P4), el recompte d'errors per sensor (P2) i la detecció de sensor
        mort (P7).
  \item \textbf{Repartiment de la càrrega:} habitacions diferents es
        reparteixen entre particions i, per tant, entre les tasques de Kafka
        Streams. Si la clau fos només l'usuari, una llar amb moltes
        habitacions carregaria una sola partició.
  \item \textbf{Cost:} les agregacions per usuari (P5, P6) necessiten
        reagrupar les dades per \texttt{<user>}. Això implica un
        \emph{repartition topic} intern, i és un cost que acceptem a canvi
        del paral·lelisme per habitació.
\end{itemize}
Els \emph{settings} tenen la clau \texttt{<user>}, perquè la configuració és
per usuari i la compactació ha de conservar l'últim valor de cada usuari.

\subsection{Convenció de topics MQTT}

\begin{center}
\begin{tabular}{@{}lll@{}}
\toprule
Broker & Topic & Contingut \\
\midrule
Llar  & \texttt{<room>/temperature} & lectura del sensor \\
Llar  & \texttt{<room>/heatpump} & ordre per al sensor \\
Llar  & \texttt{settings} & preferències de confort de l'usuari \\
Núvol & \texttt{<user>/<room>/temperature} & lectura + \texttt{user}, \texttt{room} \\
Núvol & \texttt{<user>/<room>/heatpump} & ordre del commander \\
Núvol & \texttt{<user>/settings} & preferències (reenviades pel gateway) \\
\bottomrule
\end{tabular}
\end{center}

\subsection{Emmagatzematge, Kafka Streams i commander}
\pendent{part de la Clara: docker-compose del núvol, servei d'emmagatzematge
i model de dades d'InfluxDB, esquelet de l'aplicació Kafka Streams.}

\pendent{pipelines P2--P8 i commander.}

% ---------------------------------------------------------------- Topologia
\section{Topologia de l'aplicació Kafka Streams}
\pendent{sortida de \texttt{Topology.describe()}: sub-topologies, state
stores, topics interns i nombre de tasques.}

% ---------------------------------------------------------------- Proves
\section{Proves i resultats}

\subsection{Proves unitàries}
La lògica de traducció del gateway i la d'encaminament de l'ingestor estan
separades en funcions pures (sense xarxa), i això permet provar-les sense
cap broker. Es llancen amb \texttt{python -m unittest -v} des de cada
carpeta:
\begin{itemize}
  \item \texttt{gateway/test\_gateway.py}: 12 proves (lectura vàlida
        enriquida, valor 100 reenviat, JSON incorrecte reenviat intacte,
        UTF-8 invàlid, reenviament de \emph{settings}, topics inesperats,
        ordres d'un altre usuari, validació d'ordres). Resultat: 12/12
        correctes.
  \item \texttt{cloud/ingestor/test\_ingestor.py}: 7 proves (clau
        \texttt{<user>/<room>}, lectura mal formada ingerida tal qual, topics
        desconeguts descartats, validació de \emph{settings}). Resultat: 7/7
        correctes.
\end{itemize}
Aquestes proves cobreixen la lògica, no la connexió amb els brokers ni amb
Kafka. Les proves de l'ingestor substitueixen la llibreria
\texttt{confluent-kafka} per un mòdul buit, de manera que el productor Kafka
només es prova a la prova de punta a punta.

\subsection{Prova de punta a punta}
\pendent{executar la prova amb Docker i afegir-hi les captures:
\texttt{mosquitto\_sub} al núvol i \texttt{kafka-console-consumer} sobre
\texttt{raw-readings} mostrant les claus.}

% ---------------------------------------------------------------- Execució
\section{Com executar el projecte}

\subsection*{Requisits}
Docker i Docker Compose. Per a les proves unitàries, Python 3 amb
\texttt{paho-mqtt==1.6.1}.

\subsection*{Passos (part actual)}
\begin{enumerate}
  \item Crear la xarxa compartida (un sol cop):
\begin{lstlisting}
docker network create global_net
\end{lstlisting}
  \item Arrencar el núvol de prova (des de \texttt{cloud/}):
\begin{lstlisting}
docker compose -f docker-compose.test.yml --project-name cloud up --build
\end{lstlisting}
  \item Arrencar les llars (des de \texttt{user/user/}), cadascuna amb el seu
        nom de projecte:
\begin{lstlisting}
docker compose -f home1.yml --project-name user1 up --build
docker compose -f home2.yml --project-name user2 up --build
\end{lstlisting}
  \item Comprovar que les lectures arriben al broker del núvol:
\begin{lstlisting}
docker exec -it cloud-mosquitto_cloud-1 mosquitto_sub -t "+/+/temperature" -v
\end{lstlisting}
  \item Comprovar que arriben a Kafka amb la seva clau:
\begin{lstlisting}
docker exec -it cloud-kafka-1 /opt/kafka/bin/kafka-console-consumer.sh \
  --bootstrap-server localhost:9092 --topic raw-readings \
  --from-beginning --property print.key=true
\end{lstlisting}
  \item Publicar unes preferències des de casa de l'\texttt{user1}:
\begin{lstlisting}
docker exec -it user1-mosquitto-1 mosquitto_pub \
  -t "settings" -m '{"mode": "comfort", "min": 21, "max": 24}'
\end{lstlisting}
  \item Enviar una ordre manual a una bomba de calor (simula el commander):
\begin{lstlisting}
docker exec -it cloud-mosquitto_cloud-1 mosquitto_pub \
  -t "user1/room1/heatpump" -m '{"status": 1}'
\end{lstlisting}
\end{enumerate}

\pendent{passos del núvol complet (InfluxDB, Grafana, Kafka Streams) i de
minikube.}

% ---------------------------------------------------------------- Conclusions
\section{Conclusions}
\pendent{omplir en acabar el projecte.}

\section{Temps dedicat}
\pendent{hores dedicades per cada membre.}

% ---------------------------------------------------------------- Referències
\begin{thebibliography}{9}
\bibitem{classe} E. Gabaldon. Apunts de l'assignatura: \emph{Introduction to
  Docker}, \emph{IoT Communication (MQTT)}, \emph{Streaming Processing
  (Kafka)} i \emph{NoSQL Databases}. Computació Distribuïda, UdL.
\bibitem{paho} Eclipse Paho MQTT Python Client 1.6.1.
  \url{https://pypi.org/project/paho-mqtt/1.6.1/}
\bibitem{mosquitto} Eclipse Mosquitto.
  \url{https://mosquitto.org/}
\bibitem{kafka} Apache Kafka Documentation.
  \url{https://kafka.apache.org/documentation/}
\bibitem{confluent} confluent-kafka-python.
  \url{https://docs.confluent.io/kafka-clients/python/current/overview.html}
\bibitem{compose} Docker Compose: Networking.
  \url{https://docs.docker.com/compose/how-tos/networking/}
\end{thebibliography}

\end{document}
